%% ma: L12-B02-0002
%% lop: 12
%% chuong: 1
%% bai: 2
%% dang: 12.2.1
%% loai: DS
%% muc: [TH, VD, VD, VD]
%% dap_an: "SĐSĐ"
%% nguon: "(NBV)-ĐỀ SỐ 1-MỖI NGÀY 1 ĐỀ THI 2025.pdf (D01, MỖI NGÀY 1 ĐỀ - Đề số 1), Phần II Câu 1 (THCS THPT Lê Thánh Tông - HCM 2025)"
%% kien_thuc: "Bài 2 (giá trị lớn nhất trên đoạn), Bài 1 (cực trị, đơn điệu); đạo hàm hàm lượng giác (Lớp 11 Bài 32), phương trình lượng giác cơ bản (Lớp 11 Bài 4), công thức nhân đôi (Lớp 11 Bài 2)"
%% ghi_chu: "Trùng ý tưởng với câu cùng dạng, giáo viên đã duyệt giữ (09/10/2026). Đáp án gốc S Đ S Đ đúng (kiểm bằng sympy). Ý c: f'(x)=0 tại x=-π/3 ≈ -1,05 thuộc (-2;-1) nên f không đồng biến trên cả khoảng"
%% trang_thai: chinh-thuc
\begin{ex}
Cho hàm số $f(x)=4\sin x\cos x+2x$.
\choiceTF
{Đạo hàm của hàm số đã cho là $f'(x)=4\sin2x+2$.}
{\True Hàm số $y=f(x)$ có $4$ điểm cực trị thuộc $[-\pi;\pi]$.}
{Hàm số $y=f(x)$ đồng biến trên khoảng $(-2;-1)$.}
{\True Giá trị lớn nhất của $f(x)$ trên đoạn $\left[0;\dfrac\pi2\right]$ là $\dfrac{2\pi}{3}+\sqrt3$.}
\loigiai{Ta có $f(x)=2\sin2x+2x$.
a) $f'(x)=4\cos2x+2$. Sai.
b) $f'(x)=0\Leftrightarrow\cos2x=-\dfrac12\Leftrightarrow x=\pm\dfrac\pi3+k\pi$. Trong $[-\pi;\pi]$ có $4$ nghiệm $-\dfrac{2\pi}3,-\dfrac\pi3,\dfrac\pi3,\dfrac{2\pi}3$, đều là nghiệm đơn nên $f'$ đổi dấu: có $4$ điểm cực trị. Đúng.
c) $-\dfrac\pi3\approx-1{,}05\in(-2;-1)$ và $f'$ đổi dấu qua điểm này nên $f$ không đồng biến trên $(-2;-1)$. Sai.
d) Trên $\left[0;\dfrac\pi2\right]$: $f'(x)=0\Leftrightarrow x=\dfrac\pi3$. $f(0)=0$, $f\!\left(\dfrac\pi3\right)=\sqrt3+\dfrac{2\pi}3\approx3{,}83$, $f\!\left(\dfrac\pi2\right)=\pi\approx3{,}14$. Giá trị lớn nhất là $\dfrac{2\pi}3+\sqrt3$. Đúng.}
\end{ex}
